% Part: first-order-logic
% Chapter: axiomatic-deduction
% Section: axioms-rules-quantifiers

\documentclass[../../../include/open-logic-section]{subfiles}

\begin{document}

\olfileid{fol}{axd}{qua}

\olsection{संख्यापकांसाठी स्वयंसिद्धके आणि नियम}

\begin{defn}[संख्यापकांसाठी स्वयंसिद्धके]
संख्यापकांचे नियमन करणारी \emph{स्वयंसिद्धके} म्हणजे पुढील रूपांचे सर्व
दृष्टांत होत:
\begin{align}
\ollabel{ax:q1} & \lforall[x][!B] \lif !B(t), \\
\ollabel{ax:q2} & !B(t) \lif \lexists[x][!B].
\end{align}
कोणत्याही बंद पद~$t$ साठी.
\end{defn}

\begin{defn}[संख्यापकांसाठी नियम]
\item जर $!B \lif !A(a)$ हे !!{derivation} मध्ये आधीच आलेले असेल आणि
  $a$ ची~$\Gamma$ किंवा~$!B$ मध्ये आस्थिति नसेल, तर $!B \lif
  \lforall[x][!A(x)]$ ही निगमनाची योग्य पायरी असते.
\item जर $!A(a) \lif !B$ हे !!{derivation} मध्ये आधीच आलेले असेल आणि
  $a$ ची~$\Gamma$ किंवा~$!B$ मध्ये आस्थिति नसेल, तर
  $\lexists[x][!A(x)] \lif !B$ ही निगमनाची योग्य पायरी असते.
\end{defn}

यांपैकी कोणत्याही नियमाचा संक्षेप आपण ``\QR'' असा करू.

\end{document}
